% Source: https://artofproblemsolving.com/community/c3h1914827_what_does_the_cauchyschwarz_inequality_state
% Stable identity: aops-topic-1914827-q01
(a) Let $A=\sum a_i^2$, $B=\sum b_i^2$, $C=\sum a_ib_i$. If $A=0$, all $a_i=0$ and both sides are $0$. Otherwise, for every real $t$,
\[
0\le\sum_{i=1}^{n}(a_it-b_i)^2=At^2-2Ct+B .
\]
A quadratic with positive leading coefficient that is never negative has discriminant at most $0$: $4C^2-4AB\le0$, i.e.\ $AB\ge C^2$. Equality means the discriminant is $0$, so the quadratic has a root $t$, which happens exactly when $b_i=ta_i$ for all $i$.

(b) Apply (a) with $a_i=\sqrt{x_i}$ and $b_i=\dfrac{1}{\sqrt{x_i}}$: then $a_ib_i=1$, so
\[
\left(\sum x_i\right)\left(\sum\dfrac{1}{x_i}\right)\ge(1+1+\cdots+1)^2=n^2 .
\]
For $n=2$ this is $(a+b)\left(\dfrac{1}{a}+\dfrac{1}{b}\right)\ge4$, with equality when $a=b$.

(c) Take $a=\left(\sqrt x,\sqrt y,\sqrt z\right)$ and $b=\left(\dfrac{1}{\sqrt x},\dfrac{2}{\sqrt y},\dfrac{3}{\sqrt z}\right)$. Then $a_1b_1+a_2b_2+a_3b_3=1+2+3=6$, so
\[
(x+y+z)\left(\dfrac{1}{x}+\dfrac{4}{y}+\dfrac{9}{z}\right)\ge6^2=36 .
\]
Equality needs $\dfrac{1}{\sqrt x}=t\sqrt x$, $\dfrac{2}{\sqrt y}=t\sqrt y$, $\dfrac{3}{\sqrt z}=t\sqrt z$, i.e.\ $x:y:z=1:2:3$. For $(x,y,z)=(1,2,3)$ the value is $(1+2+3)\left(1+2+3\right)=36$, so the bound is attained and the minimum is $\boxed{36}$.
